% Generated draft: requires mathematical and source review.
\documentclass[11pt]{article}
\usepackage[margin=25mm]{geometry}
\usepackage{fontspec}
\setmainfont{Latin Modern Roman}
\usepackage{amsmath,amssymb,mathtools,unicode-math}
\setmathfont{Latin Modern Math}
\usepackage{enumitem,needspace,xurl,hyperref}
\setlength{\parindent}{0pt}
\setlength{\parskip}{5pt}
\setlength{\emergencystretch}{3em}
\newcommand{\N}{\mathbb{N}}
\newcommand{\Z}{\mathbb{Z}}
\newcommand{\Q}{\mathbb{Q}}
\newcommand{\R}{\mathbb{R}}
\newcommand{\C}{\mathbb{C}}
\newcommand{\F}{\mathbb{F}}
\newcommand{\A}{\mathbb{A}}
\newcommand{\PP}{\mathbb{P}}
\newcommand{\E}{\mathbb{E}}
\newcommand{\eps}{\varepsilon}
\newcommand{\dd}{\,\mathrm{d}}
\begin{document}
\section*{UnsolvedMath ID: 701738}
\textbf{Source problem code:} catalogue entry
\subsection{Definitions and mathematical statement}
Let \(G\) be a finite group, and let \(\operatorname{Irr}(G)\) denote its irreducible complex characters. For a subgroup \(H\leq G\), write \(1_H\) for its trivial character and \(\operatorname{Ind}_H^G(1_H)=(1_H)^G\) for the induced character. A maximal subgroup means a proper subgroup maximal under inclusion. An irreducible character \(\chi\in\operatorname{Irr}(G)\) is a \(\mathcal P\)-character if there is a maximal subgroup \(H<G\) such that
\[
\bigl\langle\chi,\operatorname{Ind}_H^G(1_H)\bigr\rangle_G>0,
\qquad
\langle\alpha,\beta\rangle_G=\frac{1}{|G|}\sum_{g\in G}\alpha(g)\overline{\beta(g)}.
\]
Thus being a constituent means occurring with positive multiplicity. The character \(\chi\) is monomial if there exist a subgroup \(K\leq G\) and a linear complex character \(\lambda\) of \(K\), meaning \(\lambda(1)=1\), such that \(\chi=\operatorname{Ind}_K^G(\lambda)\).

Define \(G^{(0)}=G\) and \(G^{(n+1)}=[G^{(n)},G^{(n)}]\), where \([L,L]\) is the subgroup generated by the commutators of elements of \(L\). The group \(G\) is solvable if \(G^{(r)}=\{1\}\) for some integer \(r\geq0\).

\textbf{Conjecture (G. Qian, Kourovka Notebook Problem 20.80).} For every finite group \(G\),
\[
G\text{ is solvable}
\quad\Longleftrightarrow\quad
\text{every }\mathcal P\text{-character of }G\text{ is monomial}.
\]
The supplied source reports that the necessity direction, namely that solvability implies that every \(\mathcal P\)-character is monomial, is true, and attributes this result to G. Qian and Y. Yang (2018). Accordingly, the remaining conjectural direction asks whether every finite group whose \(\mathcal P\)-characters are all monomial must be solvable. This report of the known direction is supplied source information, not an independently verified literature assessment.
\subsection{Short English statement}
G. Qian conjectures that a finite group is solvable exactly when every irreducible character occurring in a permutation character induced from a maximal subgroup is induced from a linear character of some subgroup. The supplied source reports the implication from solvability as known; the converse remains conjectural according to that record.
\subsection{Sources}
\begin{itemize}
\item Local UnsolvedMath registry, supplied record with permanent ID \(701738\), exact source code \texttt{catalogue entry}, titled \emph{Kourovka Notebook Problem 20.80}. The registry labels the problem open; this label has not been independently verified.
\item \emph{The Kourovka Notebook (Unsolved Problems in Group Theory)}, original problem number 20.80, proposed by G. Qian in 2022. The supplied background identifies the 20th issue (2022) and the 21st issue edition (2026), \texttt{arXiv:1401.0300v48}. The supplied collection description names E. I. Khukhro and V. D. Mazurov as editors.
\item G. Qian and Y. Yang, \emph{Commun. Algebra} \textbf{46}, no. 1 (2018), 167--175. The supplied statement cites this reference for the necessity direction. No article title or link is supplied.
\end{itemize}
\end{document}
